\chapter{总结与展望}
\label{chap:conclusion}

\section{研究总结}

本文围绕小模型在执行式知识图谱问答中的状态偏离问题，研究了真实执行状态与恢复后缀的配对监督，并通过公开任务和雷达来源问答评价其效果与适用范围。研究将数据可信度、查询执行和答案支持分开检验，使最终正确率能够结合过程与来源证据解释。

在数据层面，本文形成了保留主体、部件、量值形式、条件和来源定位的雷达读法组织方式。八个制造商来源组构成的主体材料包含 502 条读法和 96 道中文问题，保留共享主体、重复记录及有歧义内容的边界；问题、参考与模型实际证据分别保存。该工作为可查询和可复核的领域实验提供了基础，但哈希、位置与程序检查只证明相应的一致性，不能独立确认事实。当前参考与语义审阅仍属于 AI 证据，24 题固定风险分层核验材料虽已准备，真人核验完成数仍为零。

在方法层面，本文从初始模型的真实失败轨迹中构造 418 个恢复配对。恢复组在共同成功前缀后保留分歧动作及观察，通过句柄重映射接入可执行参考后缀，并仅监督后续模型动作。普通组与恢复组共享问题、普通轨迹及逐记录监督预算。2,000 题项目留出评价中，两组续训种子的准确率增量分别为 4.65 和 4.85 个百分点，预定平均增量为 4.75 个百分点，按问题配对 bootstrap 的 95\% 区间为 $[3.53,5.98]$。这一结果支持恢复配对监督相对本项目普通续训基线的整体收益。

恢复组还减少了未完成题目与无效调用，累计总 token 相对对应普通组下降，但相对一次性程序生成仍约为二十四倍。质量改善因此伴随明确的交互成本。训练配对中 414 个分歧动作可以执行，只有 4 个被拒绝；详细报错文字屏蔽诊断未显示稳定的额外收益。因此，现有结果主要支持偏离状态后的恢复训练，尚不能解释为模型已经掌握错误语义。短程强化学习扩展也没有提供明确增量，未构成本文主方法贡献。

在领域应用层面，共同合成适配改善了有限查询接口的可用性，但四来源 24 题评价中，恢复组相对普通组分别少答对一题和持平，额外迁移收益尚未成立。后续八来源实验固定语义证据，以确定性条件绑定表示对比逐条平铺记录，每路分母均为 96，正文正确数由 73 提高至 80，题级绑定错误由 17 降至 11，原始联合正确数由 17 提高至 52，尚低于原文路线的 56。另行登记的事后统一引用解析使平铺与绑定的联合正确数变为 55 和 68，原文保持 56；正文与绑定错误保持原判。原始净增三十五题缩小为十三题，说明引用可用性是联合差异的重要组成部分，剩余优势也不能全部归于事实推理。该表示结果没有训练新模型，不能替代恢复训练的领域迁移验证。

\section{研究局限}

公开实验的两个续训种子共享初始模型和训练语料，并非两次从基础模型出发的独立复现；项目留出题来自官方训练数据，不代表官方隐藏测试成绩，也不能排除预训练暴露或近似结构重叠。问题级区间条件于已有检查点，不能覆盖全部训练随机性。当前对照评价了状态插入、句柄变化与后缀监督组成的整体方案，尚未分离各部分作用；匹配监督 token 也未实现等计算量比较。

领域证据受到来源组相关性、单一模型、AI 参考及审阅暴露的共同限制。结构化表示与原文输入的信息范围、长度和整理成本不同，不能以某一指标更高宣称普遍或同成本优于原文检索增强问答。来源可追溯也不等于型号参数已经得到独立认证。这些限制要求分别保留原始结果、事后分析及负结果，而不把数据、方法和应用工作合并成未经验证的统一泛化结论。

\section{后续研究方向}

后续应先补充来源与参考的真人核验，利用已固定的 24 题材料保存独立回答、原文定位和不确定性，再交叉检查争议。该样本属于风险分层核验，即使全部通过，也不能自动将全体 96 题升级为人工金标。核验发现的问题应通过新版本明确更正，不回写原始评价记录。

在方法研究上，可围绕恢复状态的覆盖与组成开展进一步对照，区分额外状态、句柄依赖和后缀监督的作用，并将共同计算预算纳入质量与成本比较。在应用研究上，应在独立登记的新来源或其他模型设置中复验条件绑定与统一引用解析，使用相同解析规则报告正文、绑定错误和引用支持，避免继续利用已见题目搜索提示。上述方向以补足当前证据缺口为目标，将执行上的可恢复性与领域回答的可核验性继续作为分别可检验的问题。
